\documentclass[12pt]{article}
\input{structure.tex}
\renewcommand{\bottomtitlespace}{3cm}

\usepackage{etoolbox}

\makeatletter
\preto{\@verbatim}{\topsep=0pt \partopsep=0pt}
\makeatother

\begin{document}

\renewcommand{\tl}{5 сек.}
\renewcommand{\ml}{1024 MB}
\problem{Задача A2. Subway}
\addauthor{Автори: Деян Хаджи-Манич, Радослав Димитров, Мартин Копчев, Емил Инджев}{2025}{06}{07}
\renewcommand{\header}{
    НАЦИОНАЛЕН ЛЕТЕН ТУРНИР ПО ИНФОРМАТИКА\\
    Пловдив, 6-8 юни 2025 г.\\
    Група A, 11-12 клас
}

След голямо преселение в североизточна България, кметът на Радко Димитриево трябва да се справи с бързото разрастване на селото.
Неговата кампания се състояла в обещания за метро в селото и сега той трябва да построи такова.

След дълго мислене кметът е определил $N$ различни станции на метрото, номерирани с числата от $0$ до $N-1$, които ще бъдат построени. Освен това, ще бъдат изкопани $M$ насочени тунела между двойки станции, номерирани с числата от $0$ до $M-1$, където $i$-тия тунел започва от станция $A_i$ и свършва в станция $B_i$ (станциите $A_i$ и $B_i$ не са непременно различни).

В селото има $D$ на брой компании за влакове за метра, номерирани с числата от $0$ до $D-1$. Всеки тунел трябва да се обслужва от точно една компания. Всяка компания иска да обсулжва множество от циклични линии, така че всяка от $N$-те станции да участва в точно една от тези циклични линии. Тук циклична линия наричаме поредици от тунели, всеки започващ от края на предния тунел, така че последния тунел да завършва в началната станция на първия тунел и всяка станция по линията е начална на точно един тунел и съща крайна на точно един тунел. С цел да има шанс това да е възможно, всяка станция има точно $D$ тунела започващи от нея и точно $D$ тунела завършващи в нея (забележете, че следва, че $M = ND$).

Търси се, ако е възможно, да разпределите тунелите между компаниите, така че горните условия да са изпълнени, или да кажете, че е невъзможно. 

\subsection{Детайли по имплементацията}

Трябва да имплементирате фунцията \verb|assign_roads|:

\begin{verbatim}
bool assign_roads(int N, int M, vector<int> A, vector<int> B)
\end{verbatim}

Тя ще бъде извикана точно веднъж.
Подават ѝ се стойности $N$, $M$ и вектори \verb|A| и \verb|B| опписващи на началата и краищата на $M$-те тунела.
След изпълнението си, функцията трябва да върне \verb|true| ако кметът може да изпълни плана си и \verb|false| ако не може.

Ако Вашата функция \verb|assign_roads| върне \verb|true| трябва преди това да е извикала помощната функция \verb|answer| точно $D$ пъти, по веднъж на компания, с параметри \verb|company_id| и \verb|roads| -- компанията, за която отговаряте, и списък от тунелите, които сте решили тя да обслужва:

\begin{verbatim}
void answer(int company_id, vector<int> roads)
\end{verbatim}

Вашата програма ще се компилира заедно с грейдър предоставен от журито. Тя не трябва да имплементира \verb|main| функция или да чете от стандартния вход, или да пише към стандартния изход. Трябва да включва хедър файла \verb|subway.h|. Предоставени са Ви локални грейдър и хедър файлове, които можете да използвате, за да тествате програмата си локално.

\subsection{Ограничения}
\begin{itemize}
\item $1 \leq N \leq 30000$
\item $1 \leq D \leq 30000$
\item $1 \leq M \leq 1210000$
\item $0 \leq A_i, B_i \leq N - 1$ (забележете, че е възможно $A_i = B_i$)
\end{itemize}

\subsection{Подзадачи}
\begin{table}[H]
    \begin{tblr}{|Q[c,m]|Q[c,m]|Q[c,m]|Q[c,m]|Q[c,m]|}
        \hline
        \textbf{Подзадача} & \textbf{Точки} & $N$ & $D$ & $M$ \\
        \hline
        $1$ & $8$ & $\leq 5$ & $\leq 4 $ & $\leq 16$ \\
        \hline
        $2$ & $9$ & $\leq 5$ & $\leq 5000$ & $\leq 25000$ \\
        \hline
        $3$ & $16$ & $\leq 30000$ & $=2$ & $\leq 60000$ \\
        \hline
        $4$ & $21$ & $\leq 30000$ & $=2^K$ & $\leq 1210000$ \\
        \hline
        $5$ & $13$ & $\leq 350$ & $\leq 350$ & $\leq 2000$ \\
        \hline
        $6$ & $10$ & $\leq 350$ & $\leq 350$ & $\leq 122500$ \\
        \hline
        $7$ & $5$ & $\leq 4000$ & $\leq 70$ & $\leq 280000$ \\
        \hline
        $8$ & $6$ & $\leq 500$ & $\leq 1500$ & $\leq 600000$ \\
        \hline
        $9$ & $6$ & $\leq 1100$ & $\leq 1100$ & $\leq 1210000$ \\
        \hline
        $10$ & $6$ & $\leq 30000$ & $\leq 30000$ & $\leq 1210000$ \\
        \hline
    \end{tblr}
    \caption*{Точките за дадена подзадача се получават само ако се преминат успешно всички тестове в нея и другите подзадачи включени в нея.

    В ограничението на подзадача $4$, $K$ е произволно положително число (такова че $2^K \leq 30000$).}
\end{table}


\subsection{Пример}

В примера $N=4$, $D=3$, $M=12$ и ребрата са: \\
$(0, 1), (1, 2), (2, 0), (3, 3), (0, 1), (1, 0), (2, 3), (3, 2), (0, 2), (3, 1), (1, 0), (2, 3)$.

Тук журито ще извика:

\begin{verbatim}
assign_roads(
    4, 12,
    {0, 1, 2, 3, 0, 1, 2, 3, 0, 3, 1, 2},
    {1, 2, 0, 3, 1, 0, 3, 2, 2, 1, 0, 3})
\end{verbatim}

В случая е възможно решение, т.е. функцията трябва да върне \verb|true|, но преди това трябва да извика \verb|answer| $3$ пъти. Примерни валидни извиквания са следните:

\begin{verbatim}
answer(0, {0, 1, 2, 3})
answer(1, {4, 5, 6, 7})
answer(2, {8, 9, 10, 11})
\end{verbatim}

\end{document}
